%%%%%%%%%%%%%%%%%%%%%%%%%%%%%%%%%%%%%%%%%
%  My documentation report
%  Objetive: Explain what I did and how, so someone can continue with the investigation
%
% Important note:
% Chapter heading images should have a 2:1 width:height ratio,
% e.g. 920px width and 460px height.
%
%%%%%%%%%%%%%%%%%%%%%%%%%%%%%%%%%%%%%%%%%

%----------------------------------------------------------------------------------------
%	PACKAGES AND OTHER DOCUMENT CONFIGURATIONS
%----------------------------------------------------------------------------------------

\documentclass[11pt,fleqn]{book} % Default font size and left-justified equations

\usepackage[top=3cm,bottom=3cm,left=3.2cm,right=3.2cm,headsep=10pt,letterpaper]{geometry} % Page margins
\usepackage[dvipsnames]{xcolor}
\usepackage{lipsum} % Required for specifying colors by name

% Font Settings
\usepackage{avant} % Use the Avantgarde font for headings
%\usepackage{times} % Use the Times font for headings
\usepackage{mathptmx} % Use the Adobe Times Roman as the default text font together with math symbols from the Symbol, Chancery and Computer Modern fonts

\usepackage{microtype} % Slightly tweak font spacing for aesthetics
\usepackage[utf8]{inputenc} % Required for including letters with accents
\usepackage[T1]{fontenc} % Use 8-bit encoding that has 256 glyphs

%-------- LANGUAGE ------
\usepackage{ifthen}

\provideboolean{lfrench}\setboolean{lfrench}{false} 
\provideboolean{lenglish}\setboolean{lenglish}{false} 

%\setboolean{lfrench}{true} %FRANCAIS
\setboolean{lenglish}{true} % ENGLISH

\ifthenelse{\boolean{lfrench}}{
\usepackage[french]{varioref}
\usepackage[french]{babel}
}{}

\ifthenelse{\boolean{lenglish}}{
\usepackage[english]{varioref}
%\usepackage[english]{babel}
}{}

% ------------------
\usepackage{wrapfig}
\usepackage{listings}
\usepackage{multicol}
\usepackage{multirow}
\usepackage{colortbl}
\usepackage{booktabs}
\newcommand{\tabitem}{~~\llap{\textbullet}~~}
\usepackage{hyperref}
\usepackage{qrcode}
\usepackage{soul}
\usepackage{tabto}
\usepackage{multienum}

\definecolor{deepblue}{rgb}{0.0,0.0,0.5}
\definecolor{deepred}{rgb}{0.6,0,0}
\definecolor{deepgreen}{rgb}{0,0.5,0}
\definecolor{ocre}{RGB}{51,102,0} 
\definecolor{lightgray}{RGB}{229,229,229} 
\definecolor{palerod}{RGB}{238,232,170}
\definecolor{verttelecom}{RGB}{171,180,0}

\newcommand\pythonstyle{\lstset{
language=Python,
basicstyle=\ttfamily\footnotesize,
morekeywords={self},              % Add keywords here
frame=tb,                         % Any extra options here
showstringspaces=false
}}

\definecolor{backcolour}{rgb}{0.95,0.95,0.92}

\newcommand\termctyle{\lstset{
frame=tb,                         % Any extra options here
showstringspaces=false
}}


% Python environment
\lstnewenvironment{python}[1][]
{
\pythonstyle
\lstset{#1}
}
{}

\lstnewenvironment{termc}[1][]
{
\lstset{#1}
}
{}


% RFC
\newcommand\rfc[1]{\href{https://www.rfc-editor.org/info/rfc#1}{\textcolor{blue}{RFC #1}\index{RFC #1}}}
\newcommand\pfunction[2]{\texttt{#2}\index{Module Python!#1!#2}\reversemarginpar\marginnote{\scriptsize\textcolor{green!50!black}{#2}}\normalmarginpar\marginnote{\scriptsize\textcolor{green!50!black}{#1}}}

% boot.py

\newcommand\glos[1]{\gls{#1}\index{#1}}
\newcommand\pprog[1]{\href{https://github.com/ltn22/The-book-of-SCHC/blob/main/Programs/#1}{\texttt{#1}}\index{Programmes!#1}\marginpar{\scriptsize\textcolor{palerod!50!black}{#1}}}

% QUESTION

\usepackage[most]{tcolorbox}

\provideboolean{Response}\setboolean{Response}{true}

\newcommand{\Correct}[1]{\ifthenelse{\boolean{Response}}{#1}{\textbf{#1}}}
\newcommand{\Wrong}[1]{\ifthenelse{\boolean{Response}}{#1}{\textcolor{black!20}{#1}}}

\newwrite\tempfile
\immediate\openout\tempfile=questions.tex


\newtcbtheorem[auto counter,number within=section]{theo}%
  {Question}{fonttitle=\bfseries\upshape, 
     arc=0mm, colback=blue!5!white,colframe=blue!75!black}{Question}
     
\newcommand\Question[3]{
\begin{theo}{#1}{summation}
#2
\immediate\write\tempfile{\noexpand\textbf{Question \thetcbcounter {} page \thepage} {} }
\immediate\write\tempfile{\unexpanded{#2}\noexpand\vspace{1em}\noexpand\newline}
\immediate\write\tempfile{\unexpanded{#3}\noexpand\newline\noexpand\newline}

\end{theo}
}

% MATHS PACKAGE
\usepackage{amsmath,tikz,tkz-euclide}
\usetikzlibrary{matrix}

\newcommand*{\horzbar}{\rule[0.05ex]{2.5ex}{0.5pt}}
\usepackage{calc}

% VERBATIM PACKAGE
\usepackage{verbatim}

\usepackage{tikz}
\usetikzlibrary{scopes}
\usetikzlibrary{automata}
\usetikzlibrary[shadows]
\usetikzlibrary{shapes}
\usetikzlibrary[decorations.footprints] 
\usetikzlibrary{decorations.pathmorphing}
\usetikzlibrary{decorations.pathreplacing}
\usetikzlibrary{decorations.text}
\usetikzlibrary {arrows}
\usetikzlibrary{patterns}
\usetikzlibrary{calc}
\usetikzlibrary{external}

\usepackage{tikz-timing}

% Acronyms

\usepackage{makeidx}
\makeindex

\usepackage{acronym}

\let\oldac\ac
\renewcommand*{\ac}[1]{\oldac{#1}\index{#1}}

\usepackage{marginnote}
\setlength{\marginparwidth}{3cm}
\usepackage[a4paper,...,marginparwidth=3cm,...]{geometry}
\let\oldmarginpar\marginpar
\renewcommand\marginpar[1]{\-\oldmarginpar[\raggedleft\footnotesize #1]%
{\raggedright\footnotesize #1}}

%\setmarginnotefont{\small\itshape\color{blue}}
\newcommand\Index[1]{#1\index{#1}\reversemarginpar\marginnote{\textcolor{blue}{\scriptsize{#1}}}\normalmarginpar}
\newcommand\Indexr[1]{#1\index{#1}\marginnote{\textcolor{blue}{\scriptsize{#1}}}}
\newcommand\Indexs[1]{#1\index{#1}}
\newcommand\marge[1]{\reversemarginpar\marginnote{\textcolor{blue}{\scriptsize{#1}}}\normalmarginpar}

% Bibliography
\usepackage[style=alphabetic,sorting=nyt,sortcites=true,autopunct=true,babel=hyphen,hyperref=true,abbreviate=false,backref=true,backend=biber]{biblatex}
\addbibresource{bibliography.bib} % BibTeX bibliography file
\defbibheading{bibempty}{}

%----------------------------------------------------------------------------------------
%	VARIOUS REQUIRED PACKAGES
%----------------------------------------------------------------------------------------

\usepackage{titlesec} % Allows customization of titles

\usepackage{graphicx} % Required for including pictures
\graphicspath{{Pictures/}} % Specifies the directory where pictures are stored

\usepackage{lipsum} % Inserts dummy text

\usepackage{tikz} % Required for drawing custom shapes

\usepackage[french]{babel} % English language/hyphenation

\usepackage{enumitem} % Customize lists
\setlist{nolistsep} % Reduce spacing between bullet points and numbered lists

\usepackage{booktabs} % Required for nicer horizontal rules in tables

\usepackage{eso-pic} % Required for specifying an image background in the title page

%----------------------------------------------------------------------------------------
%	MAIN TABLE OF CONTENTS
%----------------------------------------------------------------------------------------

\usepackage{titletoc} % Required for manipulating the table of contents

\contentsmargin{0cm} % Removes the default margin
% Chapter text styling
\titlecontents{chapter}[1.25cm] % Indentation
{\addvspace{15pt}\large\sffamily\bfseries} % Spacing and font options for chapters
{\color{ocre!60}\contentslabel[\Large\thecontentslabel]{1.25cm}\color{ocre}} % Chapter number
{}  
{\color{ocre!60}\normalsize\sffamily\bfseries\;\titlerule*[.5pc]{.}\;\thecontentspage} % Page number
% Section text styling
\titlecontents{section}[1.25cm] % Indentation
{\addvspace{5pt}\sffamily\bfseries} % Spacing and font options for sections
{\contentslabel[\thecontentslabel]{1.25cm}} % Section number
{}
{\sffamily\hfill\color{black}\thecontentspage} % Page number
[]
% Subsection text styling
\titlecontents{subsection}[1.25cm] % Indentation
{\addvspace{1pt}\sffamily\small} % Spacing and font options for subsections
{\contentslabel[\thecontentslabel]{1.25cm}} % Subsection number
{}
{\sffamily\;\titlerule*[.5pc]{.}\;\thecontentspage} % Page number
[] 

%----------------------------------------------------------------------------------------
%	MINI TABLE OF CONTENTS IN CHAPTER HEADS
%----------------------------------------------------------------------------------------

% Section text styling
\titlecontents{lsection}[0em] % Indendating
{\footnotesize\sffamily} % Font settings
{}
{}
{}

% Subsection text styling
\titlecontents{lsubsection}[.5em] % Indentation
{\normalfont\footnotesize\sffamily} % Font settings
{}
{}
{}
 
%----------------------------------------------------------------------------------------
%	PAGE HEADERS
%----------------------------------------------------------------------------------------

\usepackage{fancyhdr} % Required for header and footer configuration

\pagestyle{fancy}
\renewcommand{\chaptermark}[1]{\markboth{\sffamily\normalsize\bfseries\chaptername\ \thechapter.\ #1}{}} % Chapter text font settings
\renewcommand{\sectionmark}[1]{\markright{\sffamily\normalsize\thesection\hspace{5pt}#1}{}} % Section text font settings
\fancyhf{} \fancyhead[LE,RO]{\sffamily\normalsize\thepage} % Font setting for the page number in the header
\fancyhead[LO]{\rightmark} % Print the nearest section name on the left side of odd pages
\fancyhead[RE]{\leftmark} % Print the current chapter name on the right side of even pages
\renewcommand{\headrulewidth}{0.5pt} % Width of the rule under the header
\addtolength{\headheight}{2.5pt} % Increase the spacing around the header slightly
\renewcommand{\footrulewidth}{0pt} % Removes the rule in the footer
\fancypagestyle{plain}{\fancyhead{}\renewcommand{\headrulewidth}{0pt}} % Style for when a plain pagestyle is specified

% Removes the header from odd empty pages at the end of chapters
\makeatletter
\renewcommand{\cleardoublepage}{
\clearpage\ifodd\c@page\else
\hbox{}
\vspace*{\fill}
\thispagestyle{empty}
\newpage
\fi}

%----------------------------------------------------------------------------------------
%	THEOREM STYLES
%----------------------------------------------------------------------------------------

\usepackage{amsmath,amsfonts,amssymb,amsthm} % For math equations, theorems, symbols, etc

\newcommand{\intoo}[2]{\mathopen{]}#1\,;#2\mathclose{[}}
\newcommand{\ud}{\mathop{\mathrm{{}d}}\mathopen{}}
\newcommand{\intff}[2]{\mathopen{[}#1\,;#2\mathclose{]}}
\newtheorem{notation}{Notation}[chapter]

%%%%%%%%%%%%%%%%%%%%%%%%%%%%%%%%%%%%%%%%%%%%%%%%%%%%%%%%%%%%%%%%%%%%%%%%%%%
%%%%%%%%%%%%%%%%%%%% dedicated to boxed/framed environements %%%%%%%%%%%%%%
%%%%%%%%%%%%%%%%%%%%%%%%%%%%%%%%%%%%%%%%%%%%%%%%%%%%%%%%%%%%%%%%%%%%%%%%%%%
\newtheoremstyle{ocrenumbox}% % Theorem style name
{0pt}% Space above
{0pt}% Space below
{\normalfont}% % Body font
{}% Indent amount
{\small\bf\sffamily\color{ocre}}% % Theorem head font
{\;}% Punctuation after theorem head
{0.25em}% Space after theorem head
{\small\sffamily\color{ocre}\thmname{#1}\nobreakspace\thmnumber{\@ifnotempty{#1}{}\@upn{#2}}% Theorem text (e.g. Theorem 2.1)
\thmnote{\nobreakspace\the\thm@notefont\sffamily\bfseries\color{black}---\nobreakspace#3.}} % Optional theorem note
\renewcommand{\qedsymbol}{$\blacksquare$}% Optional qed square

\newtheoremstyle{blacknumex}% Theorem style name
{5pt}% Space above
{5pt}% Space below
{\normalfont}% Body font
{} % Indent amount
{\small\bf\sffamily}% Theorem head font
{\;}% Punctuation after theorem head
{0.25em}% Space after theorem head
{\small\sffamily{\tiny\ensuremath{\blacksquare}}\nobreakspace\thmname{#1}\nobreakspace\thmnumber{\@ifnotempty{#1}{}\@upn{#2}}% Theorem text (e.g. Theorem 2.1)
\thmnote{\nobreakspace\the\thm@notefont\sffamily\bfseries---\nobreakspace#3.}}% Optional theorem note

\newtheoremstyle{blacknumbox} % Theorem style name
{0pt}% Space above
{0pt}% Space below
{\normalfont}% Body font
{}% Indent amount
{\small\bf\sffamily}% Theorem head font
{\;}% Punctuation after theorem head
{0.25em}% Space after theorem head
{\small\sffamily\thmname{#1}\nobreakspace\thmnumber{\@ifnotempty{#1}{}\@upn{#2}}% Theorem text (e.g. Theorem 2.1)
\thmnote{\nobreakspace\the\thm@notefont\sffamily\bfseries---\nobreakspace#3.}}% Optional theorem note

%%%%%%%%%%%%%%%%%%%%%%%%%%%%%%%%%%%%%%%%%%%%%%%%%%%%%%%%%%%%%%%%%%%%%%%%%%%
%%%%%%%%%%%%% dedicated to non-boxed/non-framed environements %%%%%%%%%%%%%
%%%%%%%%%%%%%%%%%%%%%%%%%%%%%%%%%%%%%%%%%%%%%%%%%%%%%%%%%%%%%%%%%%%%%%%%%%%
\newtheoremstyle{ocrenum}% % Theorem style name
{5pt}% Space above
{5pt}% Space below
{\normalfont}% % Body font
{}% Indent amount
{\small\bf\sffamily\color{ocre}}% % Theorem head font
{\;}% Punctuation after theorem head
{0.25em}% Space after theorem head
{\small\sffamily\color{ocre}\thmname{#1}\nobreakspace\thmnumber{\@ifnotempty{#1}{}\@upn{#2}}% Theorem text (e.g. Theorem 2.1)
\thmnote{\nobreakspace\the\thm@notefont\sffamily\bfseries\color{black}---\nobreakspace#3.}} % Optional theorem note
\renewcommand{\qedsymbol}{$\blacksquare$}% Optional qed square
\makeatother

% Defines the theorem text style for each type of theorem to one of the three styles above
\newcounter{dummy} 
\numberwithin{dummy}{section}
\theoremstyle{ocrenumbox}
\newtheorem{theoremeT}[dummy]{Theorem}
\newtheorem{problem}{Problem}[chapter]
\newtheorem{exerciseT}{Exercise}[chapter]
\theoremstyle{blacknumex}
\newtheorem{exampleT}{Example}[chapter]
\theoremstyle{blacknumbox}
\newtheorem{vocabulary}{Vocabulary}[chapter]
\newtheorem{definitionT}{Definition}[section]
\newtheorem{corollaryT}[dummy]{Corollary}
\theoremstyle{ocrenum}
\newtheorem{proposition}[dummy]{Proposition}

%----------------------------------------------------------------------------------------
%	DEFINITION OF COLORED BOXES
%----------------------------------------------------------------------------------------

\RequirePackage[framemethod=default]{mdframed} % Required for creating the theorem, definition, exercise and corollary boxes

% Theorem box
\newmdenv[skipabove=7pt,
skipbelow=7pt,
backgroundcolor=black!5,
linecolor=ocre,
innerleftmargin=5pt,
innerrightmargin=5pt,
innertopmargin=5pt,
leftmargin=0cm,
rightmargin=0cm,
innerbottommargin=5pt]{tBox}

% Exercise box	  
\newmdenv[skipabove=7pt,
skipbelow=7pt,
rightline=false,
leftline=true,
topline=false,
bottomline=false,
backgroundcolor=ocre!10,
linecolor=ocre,
innerleftmargin=5pt,
innerrightmargin=5pt,
innertopmargin=5pt,
innerbottommargin=5pt,
leftmargin=0cm,
rightmargin=0cm,
linewidth=4pt]{eBox}	

% Definition box
\newmdenv[skipabove=7pt,
skipbelow=7pt,
rightline=false,
leftline=true,
topline=false,
bottomline=false,
linecolor=ocre,
innerleftmargin=5pt,
innerrightmargin=5pt,
innertopmargin=0pt,
leftmargin=0cm,
rightmargin=0cm,
linewidth=4pt,
innerbottommargin=0pt]{dBox}	

% Corollary box
\newmdenv[skipabove=7pt,
skipbelow=7pt,
rightline=false,
leftline=true,
topline=false,
bottomline=false,
linecolor=gray,
backgroundcolor=black!5,
innerleftmargin=5pt,
innerrightmargin=5pt,
innertopmargin=5pt,
leftmargin=0cm,
rightmargin=0cm,
linewidth=4pt,
innerbottommargin=5pt]{cBox}

% Creates an environment for each type of theorem and assigns it a theorem text style from the "Theorem Styles" section above and a colored box from above
\newenvironment{theorem}{\begin{tBox}\begin{theoremeT}}{\end{theoremeT}\end{tBox}}
\newenvironment{exercise}{\begin{eBox}\begin{exerciseT}}{\hfill{\color{ocre}\tiny\ensuremath{\blacksquare}}\end{exerciseT}\end{eBox}}				  
\newenvironment{definition}{\begin{dBox}\begin{definitionT}}{\end{definitionT}\end{dBox}}	
\newenvironment{example}{\begin{exampleT}}{\hfill{\tiny\ensuremath{\blacksquare}}\end{exampleT}}		
\newenvironment{corollary}{\begin{cBox}\begin{corollaryT}}{\end{corollaryT}\end{cBox}}	

%----------------------------------------------------------------------------------------
%	REMARK ENVIRONMENT
%----------------------------------------------------------------------------------------

\newenvironment{remark}{\par\vspace{10pt}\small % Vertical white space above the remark and smaller font size
\begin{list}{}{
\leftmargin=35pt % Indentation on the left
\rightmargin=25pt}\item\ignorespaces % Indentation on the right
\makebox[-2.5pt]{\begin{tikzpicture}[overlay]
\node[draw=ocre!60,line width=1pt,circle,fill=ocre!25,font=\sffamily\bfseries,inner sep=2pt,outer sep=0pt] at (-15pt,0pt){\textcolor{ocre}{R}};\end{tikzpicture}} % Orange R in a circle
\advance\baselineskip -1pt}{\end{list}\vskip5pt} % Tighter line spacing and white space after remark

%----------------------------------------------------------------------------------------
%	SECTION NUMBERING IN THE MARGIN
%----------------------------------------------------------------------------------------

\makeatletter
\renewcommand{\@seccntformat}[1]{\llap{\textcolor{ocre}{\csname the#1\endcsname}\hspace{1em}}}                    
\renewcommand{\section}{\@startsection{section}{1}{\z@}
{-4ex \@plus -1ex \@minus -.4ex}
{1ex \@plus.2ex }
{\normalfont\large\sffamily\bfseries}}
\renewcommand{\subsection}{\@startsection {subsection}{2}{\z@}
{-3ex \@plus -0.1ex \@minus -.4ex}
{0.5ex \@plus.2ex }
{\normalfont\sffamily\bfseries}}
\renewcommand{\subsubsection}{\@startsection {subsubsection}{3}{\z@}
{-2ex \@plus -0.1ex \@minus -.2ex}
{.2ex \@plus.2ex }
{\normalfont\small\sffamily\bfseries}}                        
\renewcommand\paragraph{\@startsection{paragraph}{4}{\z@}
{-2ex \@plus-.2ex \@minus .2ex}
{.1ex}
{\normalfont\small\sffamily\bfseries}}

%----------------------------------------------------------------------------------------
%	HYPERLINKS IN THE DOCUMENTS
%----------------------------------------------------------------------------------------

% For an unclear reason, the package should be loaded now and not later
\usepackage{hyperref}
\hypersetup{hidelinks,backref=true,pagebackref=true,hyperindex=true,colorlinks=false,breaklinks=true,urlcolor= ocre,bookmarks=true,bookmarksopen=false,pdftitle={Title},pdfauthor={Author}}

%----------------------------------------------------------------------------------------
%	CHAPTER HEADINGS
%----------------------------------------------------------------------------------------

% The set-up below should be (sadly) manually adapted to the overall margin page septup controlled by the geometry package loaded in the main.tex document. It is possible to implement below the dimensions used in the goemetry package (top,bottom,left,right)... TO BE DONE

\newcommand{\thechapterimage}{}
\newcommand{\chapterimage}[1]{\renewcommand{\thechapterimage}{#1}}

% Numbered chapters with mini tableofcontents
\def\thechapter{\arabic{chapter}}
\def\@makechapterhead#1{
\thispagestyle{empty}
{\centering \normalfont\sffamily
\ifnum \c@secnumdepth >\m@ne
\if@mainmatter
\startcontents
\begin{tikzpicture}[remember picture,overlay]
\node at (current page.north west)
{\begin{tikzpicture}[remember picture,overlay]
\node[anchor=north west,inner sep=0pt] at (0,0) {\includegraphics[width=\paperwidth]{\thechapterimage}};
%%%%%%%%%%%%%%%%%%%%%%%%%%%%%%%%%%%%%%%%%%%%%%%%%%%%%%%%%%%%%%%%%%%%%%%%%%%%%%%%%%%%%
% Commenting the 3 lines below removes the small contents box in the chapter heading
%\fill[color=ocre!10!white,opacity=.6] (1cm,0) rectangle (8cm,-7cm);
%\node[anchor=north west] at (1.1cm,.35cm) {\parbox[t][8cm][t]{6.5cm}{\huge\bfseries\flushleft \printcontents{l}{1}{\setcounter{tocdepth}{2}}}};
\draw[anchor=west] (5cm,-9cm) node [rounded corners=20pt,fill=ocre!10!white,text opacity=1,draw=ocre,draw opacity=1,line width=1.5pt,fill opacity=.6,inner sep=12pt]{\huge\sffamily\bfseries\textcolor{black}{\thechapter. #1\strut\makebox[22cm]{}}};
%%%%%%%%%%%%%%%%%%%%%%%%%%%%%%%%%%%%%%%%%%%%%%%%%%%%%%%%%%%%%%%%%%%%%%%%%%%%%%%%%%%%%
\end{tikzpicture}};
\end{tikzpicture}}
\par\vspace*{230\p@}
\fi
\fi}

% Unnumbered chapters without mini tableofcontents (could be added though) 
\def\@makeschapterhead#1{
\thispagestyle{empty}
{\centering \normalfont\sffamily
\ifnum \c@secnumdepth >\m@ne
\if@mainmatter
\begin{tikzpicture}[remember picture,overlay]
\node at (current page.north west)
{\begin{tikzpicture}[remember picture,overlay]
\node[anchor=north west,inner sep=0pt] at (0,0) {\includegraphics[width=\paperwidth]{\thechapterimage}};
\draw[anchor=west] (5cm,-9cm) node [rounded corners=20pt,fill=ocre!10!white,fill opacity=.6,inner sep=12pt,text opacity=1,draw=ocre,draw opacity=1,line width=1.5pt]{\huge\sffamily\bfseries\textcolor{black}{#1\strut\makebox[22cm]{}}};
\end{tikzpicture}};
\end{tikzpicture}}
\par\vspace*{230\p@}
\fi
\fi
}
\makeatother % Insert the commands.tex file which contains the majority of the structure behind the template


\newcommand\pythonlst[2][]{
\lstinputlisting[language=Python, backgroundcolor=\color{palerod},   basicstyle=\footnotesize\ttfamily,
  keywordstyle=\bfseries\color{green!40!black},
  commentstyle=\itshape\color{purple!40!black},
  identifierstyle=\color{blue},
  stringstyle=\color{orange}, caption=\detokenize{#2},
  numbers=left, numberstyle=\tiny, stepnumber=2, numbersep=5pt, frame=single, #1, label=\detokenize{#2}] {Programs/#2}\index{Programmes Python!#2}%
  }

\newcommand\pythonnxt[2][]{
\lstinputlisting[language=Python, backgroundcolor=\color{palerod},   basicstyle=\footnotesize\ttfamily,
  keywordstyle=\bfseries\color{green!40!black},
  commentstyle=\itshape\color{purple!40!black},
  identifierstyle=\color{blue},
  stringstyle=\color{orange},
  numbers=left, numberstyle=\tiny, stepnumber=2, numbersep=5pt, frame=single, #1] {Programs/#2}
  }
  
  
\newcommand\pycomlst[2][]{
\lstinputlisting[language=Python, backgroundcolor=\color{gray!10},   basicstyle=\footnotesize\ttfamily,
  keywordstyle=\bfseries\color{green!40!black},
  commentstyle=\itshape\color{purple!40!black},
  identifierstyle=\color{blue},
  stringstyle=\color{orange}, caption=\detokenize{#2},
  numbers=left, numberstyle=\tiny, stepnumber=2, numbersep=5pt, frame=single, #1] {Programs/#2}\index{Programmes micro-python!#2}
  }

\newcommand\pycomnxt[2][]{
\lstinputlisting[language=Python, backgroundcolor=\color{gray!10},   basicstyle=\footnotesize\ttfamily,
  keywordstyle=\bfseries\color{green!40!black},
  commentstyle=\itshape\color{purple!40!black},
  identifierstyle=\color{blue},
  stringstyle=\color{orange},
  numbers=left, numberstyle=\tiny, stepnumber=2, numbersep=5pt, frame=single, #1] {Programs/#2}
  }



\newcommand\Youtube[1]{\begin{tcolorbox}[colback=red!5,colframe=red!75!black,title=Youtube, width=3cm]\href{#1}{\qrcode{#1}}\end{tcolorbox}}

\newcommand\fulluri[2]{\href{#2}{#1}\footnote{\url{#2}}}
%%%%%%%

\provideboolean{allchap}\setboolean{allchap}{false}

\newcommand\Input[1]{\ifthenelse{\boolean{allchap}}{\input{#1}}{}}


%-----------



\newcommand\lgf[1]{\ifthenelse{\boolean{lfrench}}{#1}{}}
\newcommand\lge[1]{\ifthenelse{\boolean{lenglish}}{#1}{}}

\newcommand\Vrai[0]{\ifthenelse{\boolean{lfrench}}{Vrai}{}\ifthenelse{\boolean{lenglish}}{True}{}}
\newcommand\Faux[0]{\ifthenelse{\boolean{lfrench}}{Faux}{}\ifthenelse{\boolean{lenglish}}{False}{}}


\ifthenelse{\boolean{lfrench}}{
\usepackage[
    type={CC},
    modifier={by-nc-nd},
    version={3.0},
]{doclicense}
} { % by default in english
\usepackage[
    type={CC},
    modifier={by-nc-nd},
    version={3.0},
    lang=English,
]{doclicense}
}


\begin{document}

\let\cleardoublepage\clearpage

%----------------------------------------------------------------------------------------
%	TITLE PAGE
%----------------------------------------------------------------------------------------

\begingroup
\thispagestyle{empty}
\AddToShipoutPicture*{\put(200,30){\includegraphics[scale=1]{Pictures/Cover.png}}} % Image background
\centering
\vspace*{5cm}
\par\normalfont\fontsize{35}{35}\sffamily\selectfont
\lgf{\textbf{LE LIVRE SUR SCHC}}
\lge{\textbf{THE BOOK OF SCHC}}
{\LARGE }\par % Book title
\vspace*{1cm}
{\Huge Paul OULTRY}\par % Author name
\endgroup

%----------------------------------------------------------------------------------------
%	COPYRIGHT PAGE
%----------------------------------------------------------------------------------------

\newpage
~\vfill
\thispagestyle{empty}

%\noindent Copyright \copyright\ 2014 Andrea Hidalgo\\ % Copyright notice

\noindent \textsc{IMT Atlantique}\\
\doclicenseThis


\noindent \textit{\lgf{Publié le \today}\lge{Published \today}} % Printing/edition date

%----------------------------------------------------------------------------------------
%	TABLE OF CONTENTS
%----------------------------------------------------------------------------------------


\chapterimage{Pictures/Banner.png} % Chapter heading image

\pagestyle{empty} % No headers

\ifthenelse{\boolean{lfrench}}{
\renewcommand\contentsname{Table des Matières}
\renewcommand{\bibname}{Bibliographie}
}{}
\ifthenelse{\boolean{lenglish}}{
\renewcommand\contentsname{Table of Contents}
\renewcommand{\bibname}{Bibliography}
}{}

\cleardoublepage
\tableofcontents% Print the table of contents itself

%\cleardoublepage % Forces the first chapter to start on an odd page so it's on the right

\pagestyle{fancy} % Print headers again

%----------------------------------------------------------------------------------------
%	CHAPTERS
%----------------------------------------------------------------------------------------
\cleardoublepage

\lgf{\chapter*{Acronymes}}
\lge{\chapter*{Acronyms}}

\begin{multicols}{2}
\begin{acronym}
%\acro{3GPP}{3rd Generation Partnership Project}
%\acro{ABP}{Authentication By Personalisation}
%\acro{ADSL}{Asymmetric Digital Subscriber Line}
%\acro{AMQP}{Advanced Message Queuing Protocol}
\acro{AS}{Application Server}
%\acro{ASCII}{American Standard Code for Information Interchange}
%\acro{BLE}{Bluetooth Low Energy}
\acro{CBOR}{Concise Binaire Object Representation}
\acro{CoAP}{Constrained Application Protocol}
%\acro{Cosem}{Companion Specification for Energy Management}
%\acro{CRC}{Cyclic Redundancy Check}
%\acro{CSV}{Comma Separated Values}
%\acro{DLMS}{Device Language Message Specification}
%\acro{DTT}{Digital Terrestrial Television}
%\acro{DR}{Data Rate}
%\acro{GSMA}{GSM Association}
%\acro{HTML}{HyperText Markup Language}
%\acro{HTTP}{HyperText Transport Protocol}
%\acro{HTTPS}{HyperText Transport Protocol Secure}
%\acro{IANA}{Internet Assigned Numbers Authority}
%\acro{IBAN}{International Bank Account Number}
%\acro{IEEE}{Institute of Electrical and Electronics Engineers}
%\acro{IETF}{Internet Engineering Task Force}
%\acro{IoT}{Internet of Things}
\acro{IP}{Internet Protocol}
\acro{IPv4}{Internet Protocol version 4}
\acro{IPv6}{Internet Protocol version 6}
%\acro{IPSO}{IP for Smart Objects}
%\acro{ITU}{International Telecommunication Union}
%\acro{IRI}{International Resource Identifier}
%\acro{ISBN}{International Standard Book Number}
%\acro{ISO}{International Standardization Organization}
%\acro{JMS}{Java Messaging Service}
%\acro{JSON}{JavaScript Object Notation}
%\acro{JSON-LD}{JavaScript Object Notation  for Linked Data}
%\acro{LCIM}{Levels of Conceptual Interoperability Model}
\acro{LPWAN}{Low Power Wide Area Network}
%\acro{LwM2M}{Lightweight Machine to Machine}
\acro{LNS}{LoRaWAN Network Server}
%\acro{MQTT}{Message Queuing Telemetry Transport}
%\acro{NAT}{Network Address Translation}
%\acro{NGW}{Network GateWay}
%\acro{NIDD}{Non IP Data Delivery}
%\acro{OMA}{Open Mobile Alliance}
%\acro{OTAA}{Over The Air Authentication}
%\acro{OVH}{On Vous Herbèrge}
%\acro{PAC}{Porting Authorization Code}
%\acro{REST}{REpresentational State Transfer}
%\acro{RFC}{Request For Comments}
%\acro{RGW}{Radio GateWay}
%\acro{RNIPP}{Répertoire National d'Identification des Personnes Physiques}
%\acro{RSSI}{Received Signal Strength Indicator}
%\acro{RTT}{Round Trip Time}
%\acro{SCEF}{Service Capability Exposure Function}
%\acro{SenML}{Sensor Measuring List}
%\acro{SCHC}{Static Context Header Compression}
%\acro{SF}{Spreading Factor}
%\acro{SNR}{Signal to Noise Ratio}
%\acro{SSID}{Service Set IDentifier}
%\acro{STIC}{Sciences et Technologies de l’Information et de la Communication}
%\acro{TCP}{Transmission Control Protocol}
%\acro{TLV}{Type Length Value}
%\acro{TNT}{Télévision Numérique Terrestre}
%\acro{TTN}{The Things Network}
%\acro{UDP}{User Datagram Protocol}
%\acro{UIT}{Union internationale des télécommunications}
%\acro{UNB}{Ultra Narrow-Band}
%\acro{URI}{Universal Resource Identifier}
%\acro{URL}{Univeral Resource Locator}
%\acro{URN}{Univeral Resource Name}
%\acro{VPS}{Virtual Private Server}
%\acro{W3C}{World Wide Web Consortium}
%\acro{WWW}{World Wide Web}
%\acro{XML}{Extensible Markup Language}
%\acro{XMPP}{Extensible Messaging Protocol et Presence}
\end{acronym}
\end{multicols}

\setboolean{allchap}{true} % true: take all, false take nothing only /input

\Input{Part01.0-Intro.tex}
\Input{Part02.0-Rules.tex}

\chapter{Compressing Ping6}

\section{Introduction}

In this chapter, we are going to setup the compression of a ping6 request. An ICMPv6 request is composed of an IPv6 header, followed by a \Index{ICMPv6 request} message.

\begin{termc}[backgroundcolor=\color{gray!10}, basicstyle=\ttfamily\tiny, escapechar=@]
Internet Protocol Version 6, Src: 2a01:cb08:903a:bd00:a8c4:5c6d:c2b5:84be, Dst: 
2001:470:1f21:1d2::1
    0110 .... = Version: 6
    .... 0000 0000 .... .... .... .... .... = Traffic Class: 0x00 (DSCP: CS0, ECN: Not-ECT)
        .... 0000 00.. .... .... .... .... .... = Differentiated Services Codepoint: Default (0)
        .... .... ..00 .... .... .... .... .... = Explicit Congestion Notification: Not ECN-Capable Transport (0)
    .... .... .... 1000 0000 1000 0000 0000 = Flow Label: 0x80800
    Payload Length: 16
    Next Header: ICMPv6 (58)
    Hop Limit: 56
    Source: 2a01:cb08:903a:bd00:a8c4:5c6d:c2b5:84be
    Destination: 2001:470:1f21:1d2::1
    
Internet Control Message Protocol v6
    Type: Echo (ping) request (128)
    Code: 0
    Checksum: 0x3982 [correct]
    [Checksum Status: Good]
    Identifier: 0x70ef
    Sequence: 699

    Data (8 bytes)
        Data: 609ab0db000aecb7
        [Length: 8]

0000  60 08 08 00 00 10 3a 38 2a 01 cb 08 90 3a bd 00   `.....:8*....:..
0010  a8 c4 5c 6d c2 b5 84 be 20 01 04 70 1f 21 01 d2   ..\m.... ..p.!..
0020  00 00 00 00 00 00 00 01 80 00 39 82 70 ef 02 bb   ..........9.p...
0030  60 9a b0 db 00 0a ec b7  
\end{termc}

The following compression rule can be applied to that traffic.

\begin{lstlisting}[basicstyle=\ttfamily\tiny, caption={rule 6/3 for ping traffic}, caption=fig-rule-ping, backgroundcolor=\color{yellow}]

{
    "DeviceID" : "udp:83.199.24.39:8888",
    "SoR" : [
	{
	    "RuleID": 6,
	    "RuleIDLength": 3,
	    "Compression": [
		{"FID": "IPV6.VER", "TV": 6, "MO": "equal", "CDA": "not-sent"},
		{"FID": "IPV6.TC",  "TV": 0, "MO": "equal", "CDA": "not-sent"},
		{"FID": "IPV6.FL",  "TV": 0, "MO": "ignore","CDA": "not-sent"},
		{"FID": "IPV6.LEN",          "MO": "ignore","CDA": "compute-length"},
		{"FID": "IPV6.NXT", "TV": 58, "MO": "equal", "CDA": "not-sent"},
		{"FID": "IPV6.HOP_LMT", "TV" : 64,"MO": "ignore","CDA": "not-sent"},
		{"FID": "IPV6.DEV_PREFIX","TV": "2001:470:1F21:1D2::/64",
                                               "MO": "equal","CDA": "not-sent"},
		{"FID": "IPV6.DEV_IID", "TV": "::1","MO": "equal","CDA": "not-sent"},
		{"FID": "IPV6.APP_PREFIX",         "MO": "ignore","CDA": "value-sent"},
		{"FID": "IPV6.APP_IID",         "MO": "ignore","CDA": "value-sent"},
		{"FID": "ICMPV6.TYPE",  "DI": "DW", "TV": 128,"MO": "equal","CDA": "not-sent"},
		{"FID": "ICMPV6.TYPE",  "DI": "UP", "TV": 129,"MO": "equal","CDA": "not-sent"},
		{"FID": "ICMPV6.CODE",  "TV": 0,  "MO": "equal","CDA": "not-sent"},
		{"FID": "ICMPV6.CKSUM", "TV": 0, "MO": "ignore","CDA": "compute-checksum"},
		{"FID": "ICMPV6.IDENT", "TV": 0,"MO": "ignore","CDA": "value-sent"},
		{"FID": "ICMPV6.SEQNO", "TV": 0,"MO": "ignore","CDA": "value-sent"}
		
	       ]
	} ] }
	
\end{lstlisting}



\section{Understanding compression rules}

A rule contains a list of fields description:
\begin{itemize}
    \item starting with a \Index{Field ID}. openSCHC uses a string structured with the protocol name (\texttt{IPV6} and \texttt{ICMPV6} in the rule above). Field ID uses keyword \texttt{\Index{FID}}.

    \item A \Index{Target Value} (\texttt{TV}) may be present and contains the value the rule expected to find in the field.

    \item A \Index{Direction Indicator} (\texttt{DI}) can be associated to a field. In the above example, \Indexr{\texttt{ICMPV6.TYPE}} is repeated twice for uplink and downlink. When the direction is not specified, openSCHC considers the field as bi-directionnal.

    \item A \Index{Matching Operator} (\texttt{MO}) specifies the comparison between the target Value  (\texttt{TV}) and the Field Value (\texttt{FV}) present in the packet to be compressed. \rfc{8724} defines 4 MO:

    \begin{itemize}
    \item \texttt{\Index{ignore}}: the FV is ignored therefore, the result of the field comparison is always true.
    \item \texttt{\Index{equal}}: the result of the comparison is true is the value contained in the \texttt{TV} is equal to the value contained if the \texttt{FV}. OpenSCHC allows integer and strings for the comparison.
    \item \texttt{\Index{MSB}}: is true only if the first bits of the \texttt{TV} are equal to the first bits of the \texttt{FV}. The comparison length in bits is specified in the rule with the keywork \texttt{MO.val}. For string comparison, openSCHC impose a  multiple of 8 bits for the length.
    \item \texttt{\Index{match-mapping}}: The \texttt{TV} contains an array of values. The comparison is true, if the \texttt{FV} appears somewhere in the list.
    \end{itemize}

    A rule is selected if all the comparison are true and the packet contains no other fields. 


    \item When a rule is selected, the \Index{Compression Decompression Actions} (\texttt{CDA}) are applied to the packet fields. \rfc{8724} defines several CDAs, among them:

    \begin{itemize}
        \item \Indexs{\texttt{not-sent}}: the \texttt{FV} is not sent, the decompressor will use the value stored in the rule to recover the value\marge{not-sent}.
        \item \Index{\texttt{value-sent}}: the value is explicitely sent as a residue. The decompressor will use the residue to recover the value.
        \item \Index{\texttt{LSB}}: the least significant bits are sent as residue. The decompressor will combine the \texttt{TV} and the residue to recover the field. LSB can only be used with a MSB MO. No argument are needed for LSB. The LSB size is defined from the field size and the MSB size.
        \item \Index{\texttt{mapping-sent}} : the index in the array is sent as residue. The decompressor will use the idenx to recover the value. This CDA can only be used with a match-mapping MO.
        \item \Index{\texttt{compute-*}}: the \texttt{FV} is not sent, the decompressor will use a specific algorithm to recover the value, such as compute the length or compute a checksum.
    \end{itemize}

\end{itemize}

\section{Compressing an ICMPv6 echo request packet}

In the above example, the \Indexr{\texttt{IPV6.VERSION}} field is not sent, the value 6 is stored is the rule. 
The \Indexr{\texttt{IPV6.NXT}} follows the same behavior, since ICMPv6 is expected, the \texttt{TV} is 58.

~

Note that \Indexr{\texttt{IPV6.FL}} and \Indexr{\texttt{IPV6.HOP\_LMT}}[0.45cm] use the \Indexs{ignore}/\Indexs{not-sent}\marge{ignore/not-sent} combination. During the rule selection, the \texttt{FL} is not taken into account, but during decompression the value stored in the rule is used. The reason is that IPv6 Flow Label (\texttt{IPV6.FL}) can be different from 0, in the capture the value is \texttt{0x80800}. 

~

For IPv6 Hop Limit  (\Indexr{\texttt{IPV6.HOP\_LMT}}) the value cannot be forecast, depends of how many routers are crossed and may vary during from one packet to the other. The rule ignores the value and generate a hop limit of 64. To be valid, this implies that the device will not forward anymore the packet. If the device acts as a router the Hop Limit field must be sent.

~

The address fields compression is more interesting, there is not source or destination address in the SCHC rule. SCHC defines the device  and the application \Index{addresses}\footnote{This will apply also for UDP port number}. The rule becomes independent of the direction. OpenSCHC splits the addresses into a 64 bit long prefix and a 64 bit long \Index{IID}. 

In our scenario the device address is not sent, since the value is known by the device. On the opposite, we allow any host to ping the device\footnote{This behavior has to be prohibited in a LPWAN network, for security reasons, since in this scenario we are using an UDP tunnel, there is no security issue.}. The destination address is fully sent to the device.

~~

At the ICMPv6 level, the type field is elided regarding the direction. The Echo value (128) is expected downlink; form the application to the device, and the Echo reply (129) on the other direction. In both cases the ICMPv6 code is expected to be 0.

~~~

The identifier (\Indexr{\texttt{ICMPv6.IDENT}}) and sequence number (\Indexr{\texttt{ICMPV6.SEQNO}}[0.45cm]) are not compressed and send as residue.

~~

The other fields \Indexr{\texttt{IPV6.LENGTH}} and \Indexr{\texttt{ICMPV6.CHECKSUM}}[0.45cm] are not sent and computed by the other side.

~~

The SCHC packet has the following format:

\begin{figure}[!ht] 
\centering 
\begin{tikzpicture}

\draw (0,0) node (b1)  [rectangle, draw, minimum height=1cm, minimum width = 0.3cm, fill=yellow] {};
\draw(b1) node {\small{1}};
\draw (b1.east) node (b2)  [right, rectangle, draw, minimum height=1cm, minimum width = 0.3cm, fill=yellow] {};
\draw(b2) node {\small{1}};
\draw (b2.east) node (b3)  [right,rectangle, draw, minimum height=1cm, minimum width = 0.3cm, fill=yellow] {};
\draw(b3) node {\small{0}};

\path(b1.north) -- +(0, 0.2) coordinate(hlineu);

\draw [|-|] (b1.west |- hlineu) -- coordinate(a) (b3.east |- hlineu) ;
\draw (a) node [above] {\tiny{Rule ID}};

\path(b1.south) -- +(0, -0.2) coordinate(hline);

\draw [<->] (b1.west |- hline) -- coordinate(a) (b3.east |- hline) ;
\draw (a) node [below] {\tiny{3 bits}};


\draw (b3.east) node (pref) [right, rectangle, draw, minimum height=1cm, minimum width=3cm, fill=purple!10]{};

\draw [white, thick] (pref.north) -- +(0.5, 0);
\draw [white, thick] (pref.north) -- +(-0.5, 0);
\draw [dotted] (pref.north) -- +(0.5, 0);
\draw [dotted] (pref.north) -- +(-0.5, 0);

\draw [white, thick] (pref.south) -- +(0.5, 0);
\draw [white, thick] (pref.south) -- +(-0.5, 0);
\draw [dotted] (pref.south) -- +(0.5, 0);
\draw [dotted] (pref.south) -- +(-0.5, 0);

\draw (pref) node {\small{app prefix}};

\path(b1.south) -- +(0, -0.2) coordinate(hline);


\draw (pref.east) node (iid) [right, rectangle, draw, minimum height=1cm, minimum width=3cm, fill=purple!20]{};

\draw [white, thick] (iid.north) -- +(0.5, 0);
\draw [white, thick] (iid.north) -- +(-0.5, 0);
\draw [dotted] (iid.north) -- +(0.5, 0);
\draw [dotted] (iid.north) -- +(-0.5, 0);

\draw [white, thick] (iid.south) -- +(0.5, 0);
\draw [white, thick] (iid.south) -- +(-0.5, 0);
\draw [dotted] (iid.south) -- +(0.5, 0);
\draw [dotted] (iid.south) -- +(-0.5, 0);

\draw (iid) node {\small{app iid}};

\path(b1.south) -- +(0, -0.2) coordinate(hline);

\draw [<->] (pref.west |- hline) -- coordinate(a) (pref.east |- hline) ;
\draw (a) node [below] {\tiny{64 bits}};

\draw [<->] (iid.west |- hline) -- coordinate(a) (iid.east |- hline) ;
\draw (a) node [below] {\tiny{64 bits}};

\draw (iid.east) node (ident) [right, rectangle, draw, minimum height=1cm, minimum width=1cm, fill=blue!10]{};
\draw (ident.text) node {\small{ident}};

\draw [<->] (ident.west |- hline) -- coordinate(a) (ident.east |- hline) ;
\draw (a) node [below] {\tiny{16 bits}};

\draw (ident.east) node (seqno) [right, rectangle, draw, minimum height=1cm, minimum width=1cm, fill=blue!10]{};
\draw (seqno.text) node {\small{seqno}};

\draw [<->] (seqno.west |- hline) -- coordinate(a) (seqno.east |- hline) ;
\draw (a) node [below] {\tiny{16 bits}};

\draw (seqno.east) node (data) [right, rectangle, draw, minimum height=1cm, minimum width=2cm]{};
\draw (data.text) node {\small{data}};

\draw [white, thick] (data.north) -- +(0.5, 0);
\draw [white, thick] (data.north) -- +(-0.5, 0);
\draw [dotted] (data.north) -- +(0.5, 0);
\draw [dotted] (data.north) -- +(-0.5, 0);

\draw [white, thick] (data.south) -- +(0.5, 0);
\draw [white, thick] (data.south) -- +(-0.5, 0);
\draw [dotted] (data.south) -- +(0.5, 0);
\draw [dotted] (data.south) -- +(-0.5, 0);

\draw [<->] (data.west |- hline) -- coordinate(a) (data.east |- hline) ;
\draw (a) node [below] {\tiny{\textit{n}.8 bits}};

\draw (data.east) node (p1)  [right, rectangle, draw, minimum height=1cm, minimum width = 0.3cm, fill=black!60] {};
\draw (p1.east) node (p2)  [right, rectangle, draw, minimum height=1cm, minimum width = 0.3cm, fill=black!60] {};
\draw (p2.east) node (p3)  [right,rectangle, draw, minimum height=1cm, minimum width = 0.3cm, fill=black!60] {};
\draw (p3.east) node (p4)  [right,rectangle, draw, minimum height=1cm, minimum width = 0.3cm, fill=black!60] {};
\draw (p4.east) node (p5)  [right,rectangle, draw, minimum height=1cm, minimum width = 0.3cm, fill=black!60] {};
\draw (p5.east) node (p6)  [right,rectangle, draw, minimum height=1cm, minimum width = 0.3cm, fill=black!60] {};

\draw [|-|] (p1.west |- hlineu) -- coordinate(a) (p6.east |- hlineu) ;
\draw (a) node [above] {\tiny{Padding}};

\draw [<->] (p1.west |- hline) -- coordinate(a) (p6.east |- hline) ;
\draw (a) node [below] {\tiny{5 bits}};

\draw [|-|] (b3.east |- hlineu) -- coordinate(a) (p1.west |- hlineu) ;
\draw (a) node [above] {\tiny{Compression Residue}};


\end{tikzpicture}
\caption{ICMPv6 Compression Residue} 
\label{fig-residue} 
\end{figure} 

The rule ID takes 3 bits, followed by the compression \Index{residue}s, the ICMPv6 payload and some padding bits. Since the Rule ID length is 3 bit long and the rest is byte aligned, 5 bits are needed to align the SCHC packet on a  L2-word (i.e. 1 byte).

\section {Compression process}

In OpenSCHC, SCHC Compresssion/Decompression (CD) and Frgamentation/Reassembly (FR) processes are done by the SCHC Machine.

\begin{figure}[!ht] 
\centering 


\begin{tikzpicture}

\draw (0,0) node (scapy) [rectangle, draw, minimum width=2cm, minimum height=1cm] {};

\draw (scapy) node {\small{scapy}};

\draw [<-] [thick] (scapy.south) -- coordinate(d) +(0, -2.5) coordinate (a);
\draw [thick] (a) -- +(6, 0) -- +(-2, 0) coordinate (b);

\draw(b) node [above right] {\tiny{Network}};

\draw [<->, double] (scapy.east) -- +(2, 0) coordinate (c);

\draw (c) node (SCHC) [right, rectangle, draw, minimum width=3cm, minimum height=2cm] {};
\draw (SCHC) node {SCHC Machine};
\draw (SCHC.south) node (L2)  [below, rectangle, draw, minimum width = 3cm, minimum height=0.5cm] {};
\draw (L2) node {Layer 2};

\draw [->] (L2.south) -- (L2.south |-a);


\end{tikzpicture}

\caption{Ping gateway architecture} 
\label{fig-schc-archi} 
\end{figure} 

Figure~\vref{fig-schc-archi} gives the architecture of the ping gateway. 
\Index{Scapy} \textit{sniff} function gets traffic on the network. 
This traffic is composed of IPv6 packets that could be compressed and tunneled SCHC packets that are decompressed. 



\subsection{The code}
\label{sec-compr_code}


The scapy module calls the SCHC Machine, in charge of CD-FR processes. Layer 2 allows to send decompressed packet or a SCHC packet on the network. 

\subsubsection{Main program}

\pythonlst[firstline=1,lastline=15, firstnumber=1]{ping_core1.py}


This program~\pprog{ping\_core1.py} shows how to compress ICMPv6 packets and send them on a UDP tunnel to a device. It starts with the modules importation from scapy, openSCHC and some regular python modules. Note the importation of the \Index{\texttt{scapy\_connection}} module, located on the working directory. This module implements the scheduler and the methods to send packets.

\pythonnxt[firstline=18,lastline=22, firstnumber=18]{ping_core1.py}

The compression process starts with the creation of a Rule Manager \texttt{rm} to include the rules from file \texttt{icmp.json} (cf. Listing~\vref{rule-icmp1}), which contains  the compression rules we saw below and a NoAck fragmentation rule.


\begin{lstlisting}[caption={rule icmp1.json}, backgroundcolor=\color{yellow}, label=rule-icmp1, basicstyle=\ttfamily\tiny],
{
    "DeviceID" : "udp:83.199.24.39:8888",
    "SoR" : [
	 {
	    "RuleID": 6,
	    "RuleIDLength": 3,
	    "Compression": [
		{"FID": "IPV6.VER", "TV": 6, "MO": "equal", "CDA": "not-sent"},
		{"FID": "IPV6.TC",  "TV": 0, "MO": "equal", "CDA": "not-sent"},
		{"FID": "IPV6.FL",  "TV": 0, "MO": "ignore","CDA": "not-sent"},
		{"FID": "IPV6.LEN",          "MO": "ignore","CDA": "compute-length"},
		{"FID": "IPV6.NXT", "TV": 58, "MO": "equal", "CDA": "not-sent"},
		{"FID": "IPV6.HOP_LMT", "TV" : 255,"MO": "ignore","CDA": "not-sent"},
		{"FID": "IPV6.DEV_PREFIX","TV": "2001:470:1F21:1D2::/64",
                                               "MO": "equal","CDA": "not-sent"},
		{"FID": "IPV6.DEV_IID", "TV": "::1","MO": "equal","CDA": "not-sent"},
		{"FID": "IPV6.APP_PREFIX",         "MO": "ignore","CDA": "value-sent"},
		{"FID": "IPV6.APP_IID",         "MO": "ignore","CDA": "value-sent"},
		{"FID": "ICMPV6.TYPE",  "DI": "DW", "TV": 128,"MO": "equal","CDA": "not-sent"},
		{"FID": "ICMPV6.TYPE",  "DI": "UP", "TV": 129,"MO": "equal","CDA": "not-sent"},
		{"FID": "ICMPV6.CODE",  "TV": 0,  "MO": "equal","CDA": "not-sent"},
		{"FID": "ICMPV6.CKSUM", "TV": 0, "MO": "ignore","CDA": "compute-checksum"},
		{"FID": "ICMPV6.IDENT", "TV": 0,"MO": "ignore","CDA": "value-sent"},
		{"FID": "ICMPV6.SEQNO", "TV": 0,"MO": "ignore","CDA": "value-sent"}
		
	    ]
	 },{
		"RuleID" : 12,
		"RuleIDLength" : 11,
		"Fragmentation" : {
			"FRMode": "NoAck",
			"FRDirection": "DW"
		}
	} 
    ]
  }
\end{lstlisting}

This rule must be adapted to your environment. The \texttt{DeviceID} must reflect the IPv4 public address of your device and the Target Value of \Indexr{\texttt{IPV6.DEV\_PREFIX}} must be set to a IPv6 prefix of your domain, such as the ones allocated by Hurricane Electric.

\pythonnxt[firstline=48,lastline=57, firstnumber=48]{ping_core1.py}

The position of the SCHC instance has to be specified. In our example, we define our role as a core. Then, we create the UDP tunnel to receive SCHC packets from devices. The default port number is \texttt{0x5C4C}.

A \texttt{lower\_layer} is created to be used by the SCHC machine when SCHC packet will have to be sent, and a \texttt{system} is created to manage events. The latter includes a \Index{scheduler}.  Both are defined in the \Index{\texttt{scapy\_connection}} module. 


From that \texttt{system} a reference to the \texttt{scheduler} is extracted to be used in the \texttt{processPkt} function as a global variable.

\pythonnxt[firstline=58,lastline=63, firstnumber=58]{ping_core1.py}

The previously created classes are regrouped arounf a \Index{SCHC machine} which will process packets for compression and fragmentation. A \texttt{schc\_machine} instance is created and the previously created rule manager is associated to it.

\pythonnxt[firstline=65,lastline=65, firstnumber=65]{ping_core1.py}


Finally, the scapy \texttt{sniff} function is called. This last line never returns. The \texttt{processPkt} function is called each time a packet is received on interface \texttt{ens3}\footnote{Scapy allowed to listen simultaneously to several interfaced, for example \texttt{["ens3", "he-ipv6"]}, but since this feature returns sometime some errors, we prefer to listen only to interface \texttt{"ens3"} which will carried tunneled IPv6 packets from Hurricane Electrics.}.

\subsubsection{Frame processing}

\pythonnxt[firstline=24,lastline=29, firstnumber=24]{ping_core1.py}

The \texttt{processPkt} function starts by calling the SCHC machine through the \pfunction{scapy\_scheduler}{run} method. This function must be called regularly, which is the case when some traffic occurs in the network, so it is important not to filter too much incoming traffic.

~

The functions looks for two types of packets:
\begin{itemize}
    \item SCHC tunneled packets coming from devices to the UDP port we defined (\texttt{0x5C4C} in our case),
    \item IPv6 tunneled packets from Hurricane Electric easily recognisable through the use of IP proto 41.
\end{itemize}

\pythonnxt[firstline=31,lastline=44, firstnumber=31]{ping_core1.py}

First, a test done to ensure that the frame has an Ethernet encapsulation\footnote{That wouldbe the case, if the Hurrican Electric interface was directly listened by scapy.}. Then we filter the frames to keep only those with:
\begin{itemize}
    \item an Ethertype equal to 0x800 indicating an IPv4 packet,
    \item an IPv4 protocol number equal to 17 indicating an UDP message,
    \item and a destination port of \texttt{0x5C4C} indicating a SCHC packet.
\end{itemize}

If all these conditions are met, we are sure that the socket has received a packet from a device, we will see in the following step how to process it.

~

Let's focus on the incoming Ping request.



\pythonnxt[firstline=45,lastline=46, firstnumber=45]{ping_core1.py}

If the IP protocol value is 41, then we have a tunneled IPv6 packet over IPv4. The first 34 bytes corresponding to the Ethernet and IPv4 headers are removed, and the resulting IPv6 packet is send to the SCHC machine for compression with the \pfunction{SCHCProtocol}{schc\_send} method. 

~~

This function will trigger the emission of SCHC packet: 
\begin{itemize}
    \item if the compressed SCHC packet enters in one frame, it is directly sent,
    \item otherwise a fragmentation rule is applied and several SCHC packets can be sent.
\end{itemize}

The function returns:
\begin{itemize}
    \item \texttt{None} if a compression rule and optionally a fragmentation rule are found.
    \item False if a rule is missing.
\end{itemize}


~~


The SCHC packet will be sent using \Index{\texttt{scapy\_connection}} functions.

 \begin{figure}[!ht] 
\centering 


	\begin{tikzpicture}[scale=0.8, transform shape]

	\draw (0,0) node(SCHC)[rectangle, minimum size=2cm, draw, fill=palerod]{}; 
	\draw(SCHC) node {\small{SCHC}};
	\draw(SCHC.north) node [above]{\small{Core}};
	
	\draw(SCHC.south) -- ++(0, -1) coordinate(a);
	\draw (a) -- ++(-2, 0);
	\draw (a) -- ++(2, 0);
	
	\path (SCHC) -- ++(5, 0) coordinate (HE);
	
	\draw (HE) node (HEbox)[cloud, minimum width=3cm, minimum height= 2cm, draw, fill=pink!20] {};
	\draw (HE) node  {\small{ping}};
	\draw (HEbox.north) node [above] {\small{Application}};

	
	\path (SCHC) -- ++(-7, 0) coordinate (dev);
	\draw (dev) node(SCHC_dev)[rectangle, minimum size=2cm, draw, fill=gray!10]{}; 
	\draw(dev) node {\small{SCHC}};
	\draw (dev) node[cloud, minimum size=4cm, draw]{}; 
	\draw (SCHC_dev.north) node [above] {\small{Device}};
	
	
	\draw [very thick] (dev.south) to [bend right=80] node [below, sloped] {UDP}  node [above, sloped] {SCHC packet} (SCHC.south);
	
		\draw [double, blue, ->] (HE) to [bend left=45] (SCHC.south);

\draw [->, decoration={text color=blue,
            text along path,
            text={Ping Request},
            text align={center},
            raise=0.2cm, reverse path}, decorate,blue] (HE) to [bend left=45]  (SCHC.south);


	\end{tikzpicture}
\caption{SCHC core processing a Ping Request from an Application} 
\label{fig-topo} 
\end{figure} 

\subsection{The execution}

Start the program, in sudo mode:

\begin{termc}[backgroundcolor=\color{palerod}, basicstyle=\ttfamily\small, escapechar=@]
$ sudo python3.9 ping_core1.py 
\end{termc}

The program will display the rules and cursor spins, indicating the SCHC machine is running\footnote{The number of received packets determines the spinning speed. It takes 10 packets to change the cursor appearance.}. 

Start pinging the device, from any host in the Internet with IPv6 connectivity. to the IPv6 address of the device defined in the rule, the \texttt{-c 1} limits the number of ping messages. 

\begin{termc}[backgroundcolor=\color{pink!20}, basicstyle=\ttfamily\small, escapechar=@]
$ ping6 2001:470:1f21:1d2::1 -c 1
\end{termc}

Of course there is no answer to the ping, the openSCHC core instance displays some messages. They can be avoided, if the \texttt{verbose} argument is set to False or not specified when creating the \texttt{schc\_machine} object. 

\begin{termc}[backgroundcolor=\color{palerod}, basicstyle=\ttfamily\tiny, escapechar=@]
schc recv-from-l3 None None b'`\x05\x0c\x00\x00\x10:8*\x01\xcb\x08\x90:\xbd\x00I\xe0\xa3\xec\x01Vv\x9c \x01\x04p...'
schc parser {('IPV6.VER', 1): [6, 4], ('IPV6.TC', 1): [0, 8], ('IPV6.FL', 1): [330752, 20]...} 
schc compression rule {'RuleID': 6, 'RuleIDLength': 3, 'Compression': [{'FID': 'IPV6.VER', 'FL': 4, ...}
schc compression result b'\xc5\x40\x39\x61\x12\x07\x57\xa0\x09\x3c\x14\x7d\x80\x2a\xce...'/227
schc fragmentation not needed size=227
\end{termc}
 
This trace gives the openSCHC compression process divided into several steps:
\begin{itemize}
\item Parse the packet; from a sequence of byte received on the network, create a list of fields containing the field identification and their associated value.
\item Find a valid compression rule; ask the rule manager to find a rule matching the parsed packet. The rule selection will also provide the device ID.
\item Apply the compression rule.
\item Send the SCHC packet to the device ID.
\end{itemize}

~~ 

In more details:
\begin{itemize}
    \item The first line (\texttt{recv-from-l3}) dumps the original IPv6 packet received by the compressor, corresponding in hexadecimal to:
\end{itemize}

\begin{termc}[backgroundcolor=\color{palerod}, basicstyle=\ttfamily\tiny, escapechar=@]
60050c0000103a382a01cb08903abd0049e0a3ec0156769c200104701f2101d20000000000000001800051fb48b20000609f882600060ed2
\end{termc}

\begin{itemize}
    \item The second line (\texttt{parser}) shows three elements returned by the parser. The first one is the list of fields, the second one is the data following the list of fields and the third one is a status code. The parsed headers are displayed figure~\ref{list-parsed-icmpv6}:
\end{itemize}

\begin{termc}[backgroundcolor=\color{palerod}, basicstyle=\ttfamily\small, escapechar=@, label=list-parsed-icmpv6, caption={Parsed IPv6/ICMPv6 header fields}]
{('ICMPV6.CKSUM', 1): [20987, 16],
 ('ICMPV6.CODE', 1): [0, 8],
 ('ICMPV6.IDENT', 1): [18610, 16],
 ('ICMPV6.SEQNO', 1): [0, 16],
 ('ICMPV6.TYPE', 1): [128, 8],
 ('IPV6.APP_IID', 1): [b'I\xe0\xa3\xec\x01Vv\x9c', 64],
 ('IPV6.APP_PREFIX', 1): [b'*\x01\xcb\x08\x90:\xbd\x00', 64],
 ('IPV6.DEV_IID', 1): [b'\x00\x00\x00\x00\x00\x00\x00\x01', 64],
 ('IPV6.DEV_PREFIX', 1): [b' \x01\x04p\x1f!\x01\xd2', 64],
 ('IPV6.FL', 1): [330752, 20],
 ('IPV6.HOP_LMT', 1): [56, 8, 'fixed'],
 ('IPV6.LEN', 1): [16, 16, 'fixed'],
 ('IPV6.NXT', 1): [58, 8, 'fixed'],
 ('IPV6.TC', 1): [0, 8],
 ('IPV6.VER', 1): [6, 4]}
\end{termc}
\begin{itemize}

As shown in figure~\vref{list-parsed-icmpv6}, a header description is a dictionnary where keys are tuple Field ID and position\footnote{In this example, position is always 1 since no field is repeated several time.}, and the value is the tuple containing field value and field size in bits.
\end{itemize}

\begin{itemize}
 The next return element is the data, i.e. the bytes following the parsed header:
\end{itemize}

\begin{termc}[backgroundcolor=\color{palerod}, basicstyle=\ttfamily\small, escapechar=@]
b'`\x9f\x88&\x00\x06\x0e\xd2'
\end{termc}
\begin{itemize}
and the error code is None since the parser recognize all the fields.
\end{itemize}

\begin{itemize}
This no surprise, the rule 6/3 matches. 
\end{itemize}


\begin{itemize}

\item The third line \texttt{compression result} gives the SCHC packet. Note the \texttt{/227}\footnote {227 \% 8 = 3.  Since the all the residues are byte alligned, the 3 represent the rule ID length. 5 bits of padding will have to be addeed.} at the end, indicating the length in bits. Converted in hexadecimal, we have:
\end{itemize}


\begin{termc}[backgroundcolor=\color{palerod}, basicstyle=\ttfamily\small, escapechar=@]
b'c5403961120757a0093c147d802aced3891640000c13f104c000c1da40'
\end{termc}


\begin{itemize}
    \item Finally, no fragmentation is required, so the SCHC packet is directly sent on the tunnel. A frame capture of UDP frame with port 0x5C4C gives:
\end{itemize}



\begin{termc}[backgroundcolor=\color{palerod}, basicstyle=\ttfamily\footnotesize]
>sudo tcpdump -nXi ens3 udp port 0x5C4C
tcpdump: verbose output suppressed, use -v or -vv for full protocol decode
listening on ens3, link-type EN10MB (Ethernet), capture size 262144 bytes
16:41:48.609729 IP 51.91.121.182.23628 > 83.199.24.39.8888: UDP, length 29
	0x0000:  4500 0039 ac95 4000 4011 751f 335b 79b6  E..9..@.@.u.3[y.
	0x0010:  53c7 1827 5c4c 22b8 0025 1936[c540 3961  S..'\L"..%.6.@9a
	0x0020:  1207 57a0 093c 147d 802a ced3 8e5f c000  ..W..<.}.*..._..
	0x0030:  0c13 fbb5 8000 d951 80]                   .......Q.

\end{termc}

\section{The Decompression Process}

 Now that we send SCHC packet to the device, lets process the SCHC packet on this side.
 
\subsection{Decompression}

Let's do the same operation on the device side. The code is almost the same, as the core SCHC. It is important to note that the rules are excactly the same as the one we used in the core SCHC.

\pycomlst[firstline=46,lastline=57, firstnumber=46]{ping_device1.py}

The position is changed to \texttt{T\_POSITION\_DEVICE}.  Since the \texttt{device\_id}, in our example, is based on a public IP address, the call to \texttt{https://api.ipify.org} returns the IPv4 public address of the device. We create a UDP socket for port 8888.


\pycomnxt[firstline=59,lastline=69, firstnumber=59]{ping_device1.py}

The \texttt{lower\_layer}, \texttt{system} and \texttt{schc\_machine} are declared the same way as in the core SCHC.  The interface name is adapted to its name on the device side.

\pycomnxt[firstline=35,lastline=41, firstnumber=35]{ping_device1.py}


It's time now, to look how a SCHC packet is processed. If the IP protocol is 17, the upper layer is UDP. We check the destination port (here 8888) to see if the message is for us. If it is the case, then we can get officially the message from the socket, to clear the buffers. We read the socket and get the message (\texttt{schc\_pkt}) and the device address (\texttt{addr}) which send the SCHC packet. 

~~

Then we call the function \pfunction{SCHCProtocol}{schc\_recv} with the device ID we computed before and the SCHC packet we just received. This function  returns:
\begin{itemize}
    \item \texttt{None}. This is the normal behavior when receiving a fragment. Only the last fragment will return the full messages. Intermediary fragments will return this value.
    \item a byte array with the uncompressed packet.
    \item \texttt{False} when a error occurs.
\end{itemize}

Since we are just doing compression, the \texttt{schc\_recv} function returns the full packet. You may notice that some fields have changed. That's the case for \Indexr{IPV6.FL}, \Indexr{PV6.HOP\_LMT}[0.4cm]. 



\subsection{Device optimization}

In the previous program, we started to reconstruct the packet. It can be then processed normally or even forwarded to another destination. Nevertheless, we can optimize the behavior and process directly the SCHC packet. This simplify the protocol stack implementation in a very constrained device. 

~~

\pycomlst[firstline=7,lastline=11, firstnumber=7]{ping_device2.py}

We have to identify the rule in the SCHC packet, we just receive. We need to transform the Byte Array, into a bit array. Therefore, the \Index{\texttt{gen\_bitarray}} module is imported\footnote{Other importations will be used when IPv6 packet will be manipulated.}.

\pycomlst[firstline=65,lastline=74, firstnumber=65]{ping_device2.py}

The SCHC packet is first transformed into a bit array with the \pfunction{gen\_bitarray}{BitBuffer} constructor. 

~~

We can call the function \pfunction{gen\_rulemanager}{FindRuleFromSCHCpacket} with the device ID and the SCHC packet. If a rule ID is found, the function will return it\footnote{We are very optimistic in this example. If no rule is found, the function will return None and the program will crash.}.

~~~

We check the rule, to see if it correspond to the ping dedicated rule 6/3. We need to test the value and the length, since two different rules may have the same value (remember chapter~\vref{chap-rules}).

~~~

If the rule in the SCHC packet corresponds, we directly answer to the core. We use a trick in the rule definition linked to direction (cf. rule~\vref{rule-icmp1}). 

\begin{lstlisting}[basicstyle=\ttfamily\tiny, caption={extract from rule 6/3 for ping traffic}, backgroundcolor=\color{yellow}]

{
    "DeviceID" : "udp:83.199.24.39:8888",
    "SoR" : [
	{
	    "RuleID": 6,
	    "RuleIDLength": 3,
	    "Compression": [
...
		{"FID": "ICMPV6.TYPE",  "DI": "DW", "TV": 128,"MO": "equal","CDA": "not-sent"},
		{"FID": "ICMPV6.TYPE",  "DI": "UP", "TV": 129,"MO": "equal","CDA": "not-sent"},
...
	       ]

\end{lstlisting}

As you have may notice, the \Indexr{ICMPV6.TYPE} is defined twice: one for downlink direction, (i.e. from the network to the device) and one in the uplink direction. For the former, the code for an ICMPv6 Echo Request (128) is set and for the latter, it corresponds to the code of an ICMPv6 Echo Reply (129).

So if we just echo the SCHC packet, it will be as a response, and application address, sequence number and identifier it carries will be included in the response. Therefore we just need to send the packet to the core SCHC we got the address when we read the socket.



We can optimize the process on the device. The rule 6/3 (cf. rule figure~\vref{rule-icmp1}) has been associated to the ping traffic. For the downlink, the type is an Echo Request (128) and for the uplink the type is an Echo Reply (129). The other fields remain unchanged in both directions.

Let's continue with this hack and see how the core should process the SCHC packet.

 \begin{figure}[!ht] 
\centering 


	\begin{tikzpicture}[scale=0.8, transform shape]

	\draw (0,0) node(SCHC)[rectangle, minimum size=2cm, draw, fill=palerod]{}; 
	\draw(SCHC) node {\small{SCHC}};
	\draw(SCHC.north) node [above]{\small{Core}};
	
	\draw(SCHC.south) -- ++(0, -1) coordinate(a);
	\draw (a) -- ++(-2, 0);
	\draw (a) -- ++(2, 0);
	
	\path (SCHC) -- ++(5, 0) coordinate (HE);
	
	\draw (HE) node (HEbox)[cloud, minimum width=3cm, minimum height= 2cm, draw, fill=pink!20] {};
	\draw (HE) node  {\small{ping}};
	\draw (HEbox.north) node [above] {\small{Application}};

	
	\path (SCHC) -- ++(-7, 0) coordinate (dev);
	\draw (dev) node(SCHC_dev)[rectangle, minimum size=2cm, draw, fill=gray!10]{}; 
	\draw(dev) node {\small{SCHC}};
	\draw (dev) node[cloud, minimum size=4cm, draw]{}; 
	\draw (SCHC_dev.north) node [above] {\small{Device}};
	

	\draw [very thick,blue, <->] (dev.south) to [bend right=80]  (SCHC.south);	
	\draw [very thick, blue, <->, decoration={text color=blue,
            text along path,
            text={rule ID+residue+data},
            text align={center},
            raise=0.2cm}, decorate] (dev.south) to [bend right=80] (SCHC.south);
	
	\draw [double, grey, ->] (HE) to [bend left=45] (SCHC.south);




	\end{tikzpicture}
\caption{SCHC core processing a Ping Request from an Application} 
\label{fig-topo} 
\end{figure} 

\section{Generating packets}

Now, the SCHC core instance receives SCHC packets from the device through the UDP tunnel. We can decompress the SCHC message and send it to the Internet.

\pythonnxt[firstline=31,lastline=44, firstnumber=31]{ping_core2.py}

The processing is almost the same as for the device, if a UDP message with destination port corresponding to the tunnel (i.e. \texttt{0x5C4C}) is detected in the packet, then the socket is read. It differs in the way the device is identified. Here, the sender address returned by the socket is used to build the \Index{device ID} (stored in \texttt{other\_end} variable). 

The call to the SCHC machine \pfunction{SCHCProtocol}{schc\_recv} method may return the uncompressed packet,in the scapy format\footnote{If the SCHC packet was a fragment, it will have been buffered until the last fragment is received. In that case None would have been returned.}. 

If the packet is returned, the \pfunction{scapy}{show} function display it and then it is sent on the Hurricane Electric interface.



\Input{Part04.0-Intro-NoAck.tex}

%%%%%%%%%%%%%%%%%%%%%%%%%%%%%%%%%%%%%



\immediate\closeout\tempfile
\setboolean{Response}{false}

\cleardoublepage
\lgf{\chapter{Réponses aux questions}}
\lge{\chapter{Answers to the questions}}
\textbf {Question 1.6.1{} page 16} {} 
A quoi correspond la clé \texttt {'u' : 'Cel'} que l'on retrouve dans la structure précédente ? \vspace {1em}\newline 
Unité = degrés Celcius\newline \newline 
\textbf {Question 1.6.2{} page 16} {} 
Dans les deux représentations JSON et CBOR, de combien la taille est-elle accrue par l'ajout des mesures effectuées ? d'où viennent ces différences ?\vspace {1em}\newline 
 Si l'on regarde le listing précédent, l'ajout des trois mesures fait augmenter la taille de 141 octets pour JSON et 84 pour CBOR. La différence vient de l'utilisation de nombre plus que de chaînes de caractères pour les clés. Ainsi 't' demande 3 caractères en JSON avec les guillemets, codé en CBOR, il faudrait 2 octets, un nombre inférieur à 23 se code sur un seul octet. Il y a egalement les virgules, espaces et fermeture de crochets qui ne sont pas présent en CBOR. Les nombres flottant comme \texttt {19.98} ont une représentation plus compacte en JSON (5 octets) qu'en CBOR où ils consomment 9 octets. Dans tous les cas l'accroissement est fortement dépendant du nom des éléments. Ici, il faut répéter à chaque fois \texttt {temperature}, \texttt {humidity} et \texttt {pressure}, soit 26 caractères. \newline \newline 
\textbf {Question 1.6.3{} page 16} {} 
Si on ne s'intéressait qu'à une seule grandeur, par exemple l'humidité. A quoi ressemblerait la structure SenML en JSON ?\vspace {1em}\newline 
 Chaque nouvelle entrée ajoute 35 octets à la structure~: [\{'bn': 'device1', 'bt': 1640110457.0, 'n': 'humidty', 'u': '\%RH', 'v': 28.46\},\\ \{'t': 10.0, 'v': 26.86\},\\ \{'t': 20.0, 'v': 26.96\},\\ \{'t': 30.0, 'v': 27.01\}]\\ \newline \newline 
\textbf {Question 1.6.4{} page 18} {} 
 Pourrait-on utiliser le champ SenML \textit {base value} pour diminuer la taille des données de pression atmosphérique~? \vspace {1em}\newline 
 Cela serait possible, si cette ressource était envoyée seule. \newline \newline 


\cleardoublepage
\printindex

\cleardoublepage
\printbibliography

\end{document}